\documentclass[hidelinks]{article}
\usepackage{stex}

\libinput{preambleCS}

\title{Formal Languages and Automata Definitions}
\author{}
\date{}

\begin{document}

\maketitle

\begin{smodule}[title={Formal Languages and Automata}]{FormalLanguagesAndAutomata}
\usemodule{ComputerScience/LexingParsingCompilingInterpreting/CompilationPhases/CompilationPhases_defs?Compilation}
\usemodule{ComputerScience/LexingParsingCompilingInterpreting/Parsing/Parsing_defs?Parsing}

    \section{Core Concepts}

    \symdecl*{formal language}
    \begin{sdefinition}[for=formal language]
        A \definame{formal language} is a set of strings over a fixed alphabet, studied independently of any particular meaning carried by the strings.
    \end{sdefinition}

    \symdecl*{language class}
    \begin{sdefinition}[for=language class]
        A \definame{language class} is a family of \sns{formal language} grouped together by a common description method, grammar restriction, or recognising machine model.
    \end{sdefinition}

    \symdecl*{automaton}
    \begin{sdefinition}[for=automaton]
        An \definame{automaton} is an abstract machine that reads input strings and changes state according to transition rules in order to recognise a \sn{formal language}.
    \end{sdefinition}

    \symdecl*{unambiguous grammar}
    \begin{sdefinition}[for=unambiguous grammar]
        An \definame{unambiguous grammar} is a \sr{Compilation?ContextFreeGrammar}{grammar} in which every string in the language has a unique \sr{Parsing?ParseTree}{parse tree}.
    \end{sdefinition}

    \section{Machine Models and Acceptance}

    \symdecl*{finite-state automaton}
    \begin{sdefinition}[for=finite-state automaton]
        A \definame{finite-state automaton} is an \sn{automaton} with finitely many states and no auxiliary memory beyond its current state.
    \end{sdefinition}

    \symdecl*{pushdown automaton}
    \begin{sdefinition}[for=pushdown automaton]
        A \definame{pushdown automaton} is an \sn{automaton} equipped with a stack, allowing it to store and later recover an unbounded amount of nested information while reading the input from left to right.
    \end{sdefinition}

    \symdecl*{nondeterministic pushdown automaton}
    \begin{sdefinition}[for=nondeterministic pushdown automaton]
        A \definame{nondeterministic pushdown automaton} is a \sn{pushdown automaton} that may have several possible transitions from the same configuration.
    \end{sdefinition}

    \symdecl*{deterministic pushdown automaton}
    \begin{sdefinition}[for=deterministic pushdown automaton]
        A \definame{deterministic pushdown automaton} is a \sn{pushdown automaton} with no choice of move in any situation: for each state, stack symbol, and input symbol (including $\lambda$), there is at most one transition, and a true-input move excludes a competing $\lambda$-move.
    \end{sdefinition}

    \symdecl*{final-state acceptance}
    \begin{sdefinition}[for=final-state acceptance]
        The \definame{final-state acceptance} condition says that a word is accepted when a complete run finishes in one of the designated accepting states.
    \end{sdefinition}

    \symdecl*{empty-stack acceptance}
    \begin{sdefinition}[for=empty-stack acceptance]
        The \definame{empty-stack acceptance} condition says that a word is accepted when a complete run of a \sn{pushdown automaton} can consume the input and empty the stack.
    \end{sdefinition}

    \symdecl*{linear-bounded automaton}
    \begin{sdefinition}[for=linear-bounded automaton]
        A \definame{linear-bounded automaton} is a Turing-machine-like automaton whose tape is restricted to a fixed finite length bounded by the input.
    \end{sdefinition}

    \symdecl*{Turing machine}
    \begin{sdefinition}[for=Turing machine]
        A \definame{Turing machine} is an automaton with a read-write head and an unbounded tape, used to model general computation.
    \end{sdefinition}

    \section{Language Classes}

    \symdecl*{regular language}
    \begin{sdefinition}[for=regular language]
        A \definame{regular language} is a \sn{formal language} recognised by a \sn{finite-state automaton} and generated by a regular grammar.
    \end{sdefinition}

    \symdecl*{context-free language}
    \begin{sdefinition}[for=context-free language]
        A \definame{context-free language} is a \sn{formal language} generated by a \sr{Compilation?ContextFreeGrammar}{context-free grammar} and recognised by a \sn{nondeterministic pushdown automaton}.
    \end{sdefinition}

    \symdecl*{Chomsky normal form}
    \begin{sdefinition}[for=Chomsky normal form]
        \definame{Chomsky normal form} is a restricted normal form for \sr{Compilation?ContextFreeGrammar}{context-free grammars} in which every production is either of the form $A \to BC$ or $A \to a$, with at most the usual special case for the empty word at the start symbol.
    \end{sdefinition}

    \symdecl*{context-sensitive language}
    \begin{sdefinition}[for=context-sensitive language]
        A \definame{context-sensitive language} is a \sn{formal language} generated by a context-sensitive grammar and recognised by a \sn{linear-bounded automaton}.
    \end{sdefinition}

    \symdecl*{recursively enumerable language}
    \begin{sdefinition}[for=recursively enumerable language]
        A \definame{recursively enumerable language} is a \sn{formal language} generated by a type-0 grammar and recognised by a \sn{Turing machine}.
    \end{sdefinition}

    \symdecl*{Chomsky hierarchy}
    \begin{sdefinition}[for=Chomsky hierarchy]
        The \definame{Chomsky hierarchy} is the proper containment chain in which every \sn{regular language} is \sn{context-free language}, every \sn{context-free language} is \sn{context-sensitive language}, and every \sn{context-sensitive language} is \sn{recursively enumerable language}.
    \end{sdefinition}

    \section{Algorithms and Properties}

    \symdecl*{CYK membership algorithm}
    \begin{sdefinition}[for=CYK membership algorithm]
        The \definame{CYK membership algorithm} is the Cocke-Younger-Kasami membership test for \sr{Compilation?ContextFreeGrammar}{context-free grammars} in \sn{Chomsky normal form}, which fills a table of non-terminals that derive substrings of the input word.
    \end{sdefinition}

    \symdecl*{closure properties for context-free languages}
    \begin{sdefinition}[for=closure properties for context-free languages]
        The \definame{closure properties for context-free languages} say that the family of \sns{context-free language} is closed under union, concatenation, and star-closure, but in general not closed under intersection or complementation.
    \end{sdefinition}

    \symdecl*{regular intersection property for context-free languages}
    \begin{sdefinition}[for=regular intersection property for context-free languages]
        The \definame{regular intersection property for context-free languages} says that if $L_1$ is a \sn{context-free language} and $L_2$ is a \sn{regular language}, then $L_1 \cap L_2$ is context-free.
    \end{sdefinition}

    \symdecl*{pumping lemma for context-free languages}
    \begin{sdefinition}[for=pumping lemma for context-free languages]
        The \definame{pumping lemma for context-free languages} says that if $L$ is context-free, then there exists an integer $n$ such that every word $w \in L$ with $|w| \geq n$ can be written as $w = uvxyz$ with $|vxy| \leq n$, $|vy| \geq 1$, and $uv^kxy^kz \in L$ for all $k \geq 0$.
    \end{sdefinition}

\end{smodule}

\end{document}
